% from ega4-13-fr.tex
\oldpage[IV-3]{198}
Or, comme $\mathcal O_y$ est supposé équidimensionnel, cette égalité s'écrit
\[
\tag{13.3.6.2}
\label{IV.13.3.6.2-fr}
  \dim_x(X_j\cap f^{-1}(y))=\dim(\mathcal O_{X,x})
  -\dim(\mathcal O_{Y,y})+\deg.\operatorname{tr}_{k(y)}k(x).
\]
Le premier membre de (13.3.6.2) est donc indépendant de $j$~; mais comme $f_j$
est équidimensionnel au point $x$, on a
$\dim_x(X_j\cap f^{-1}(y))=\dim(f_j^{-1}(y_j))$, donc le critère (13.3.3)
montre que $f$ est équidimensionnel au point $x$.

Inversement, supposons que $\mathcal O_{Y,y}$ soit un anneau universellement
caténaire et que $f$ soit équidimensionnel au point $x$~; alors (13.3.5)
$\mathcal O_{X,x}$ est équidimensionnel, et il résulte alors de (13.3.4) que
l'on a la relation
\[
  \dim(\mathcal O_{X,x})=\dim(\mathcal O_{Y,y})+e
  -\deg.\operatorname{tr}_{k(y)}k(x)
\]
où $e$ est la valeur commune des nombres
$\dim(f_j^{-1}(y_j))=\dim_x(X_j\cap f^{-1}(y))$~; mais par définition $e$ est
aussi égal à $\dim_x(f^{-1}(y))$, d'où la relation (13.3.6.1), compte tenu de
(5.6.5.2).

\begin{proposition}[13.3.7]
\label{IV.13.3.7-fr}
Soient $Y$ un préschéma, $f:X\to Y$ un morphisme localement de type fini et
équidimensionnel, $Z$ une partie fermée de $X$. Alors la fonction
\[
\tag{13.3.7.1}
\label{IV.13.3.7.1-fr}
  x\longmapsto\operatorname{codim}_x
  (Z\cap f^{-1}(f(x)),f^{-1}(f(x)))
\]
est semi-continue inférieurement dans $X$.
\end{proposition}

Comme toutes les composantes irréductibles de $f^{-1}(f(x))$ contenant $x$
ont par hypothèse la même dimension (13.3.1), on a, en vertu de (5.2.1),
\[
  \operatorname{codim}_x(Z\cap f^{-1}(f(x)),f^{-1}(f(x)))
  =\dim_x(f^{-1}(f(x)))-\dim_x(Z\cap f^{-1}(f(x))).
\]
Mais par hypothèse le premier terme du second membre est fonction
\emph{continue} de $x$ (13.3.1) et le second est fonction semi-continue
supérieurement de $x$ en vertu de (13.3.3)~; d'où la conclusion.

\begin{proposition}[13.3.8]
\label{IV.13.3.8-fr}
Soient $Y$ un préschéma, $f:X\to Y$ un morphisme localement de type fini, $x$
un point de $X$. Soient $Y'$ un préschéma, $g:Y'\to Y$ un morphisme plat,
$X'=X\times_Y Y'$, $f'=f_{(Y')}:X'\to Y'$~; si $f$ est équidimensionnel au
point $x$, alors $f'$ est équidimensionnel en tout point $x'\in X'$ au-dessus
de $x$.
\end{proposition}

La question étant locale sur $X$ et $Y$, on peut se borner au cas où toute
composante irréductible de $X$ (resp. $Y$) contient $x$ (resp. $y$)~; comme
l'image par $f$ de toute composante irréductible de $X$ est alors dense dans
une composante irréductible de $Y$ (13.3.1), on sait (2.3.5) que l'image par
$f'$ de toute composante irréductible de $X'$ est dense dans une composante
irréductible de $Y'$. D'autre part, par transitivité des fibres
(\textbf{I}, 3.6.4), il résulte de (4.2.8) et de (13.3.1) que, pour tout
$z'\in X'$, l'ensemble des dimensions des composantes irréductibles de
${f'}^{-1}(f'(z'))$ est le même que l'ensemble des dimensions des composantes
irréductibles de $f^{-1}(f(z))$, où $z$ est la projection de $z'$ dans $X$~;
d'où la conclusion en vertu de (13.3.1, $a$)).

\begin{remark}[13.3.9]
\EGAReaderTarget{ega4.c04.statement.13-3-9}{IV.13.3.9-fr}
La propriété d'équidimensionnalité d'un morphisme $f:X\to Y$ n'est \emph{pas}
stable par changement de base quelconque $g:Y'\to Y$, même lorsque $g$ est
l'injection canonique d'une composante irréductible $Y'$ de $Y$ dans $Y$
($Y'$ étant

